% =============================================================================
% 05-process-concept-pcb.tex — Lecture 5:
% Process Concept and the Process Control Block
%
% Course: Operating Systems (PCC CS-502)
% Anchor reading: Silberschatz, Galvin & Gagne, Operating System Concepts,
%                 8th ed. — Ch. 3 §3.1 (process concept, pp. 101-105: the
%                 process p.101, process state p.103, PCB p.103); page-verified
%                 via master_supporting_docs/CS502/supporting_books/SilberschatzOSC/index.md.
%                 Stallings, Operating Systems: Internals and Design
%                 Principles, 6th ed. — Ch. 3 §3.1 (what is a process, p.108),
%                 §3.2 (process states, p.111), §3.3 (process description /
%                 PCB, p.126); page-verified via .../StallingsOS/index.md.
%                 Bovet & Cesati, Understanding the Linux Kernel, 3rd ed. —
%                 Ch. 3 (processes, process descriptor, pp. 79-130);
%                 page-verified via .../LinuxKernel/index.md.
%
% Running application: one program file, ./counter, launched twice from two
% terminals — one program, two processes (P1, P2), each with its own PCB.
%
% Compile with: /compile-latex CS502/05-process-concept-pcb
% =============================================================================

\documentclass[aspectratio=169]{beamer}
\input{header}
\coursecode{CS502}

\title{Process Concept and the Process Control Block}
\subtitle{Lecture 5 --- What Is a Running Program?}
\author[Dr. Auqib Hamid Lone]{%
  \textbf{Dr. Auqib Hamid Lone}\\
  \small Assistant Professor in Computer Science \& Engineering
}
\institute[GCET Ganderbal]{%
  \small\textbf{PCC CS-502: Operating Systems}\\[0.3em]
  Department of Computer Science \& Engineering\\
  Government College of Engineering \& Technology (GCET), Ganderbal
}
\date[\today]{\small\today}

\begin{document}

\frame{\titlepage}

% =============================================================================
% ACT 0 --- ONE FILE, TWO RUNS
% =============================================================================
\section{One File, Two Runs}

\begin{frame}{Welcome: One Program, Two Terminals}
  \begin{block}{The big question}
    You compile \texttt{counter.c} into one file, \texttt{./counter}. You run
    it in two terminals at the same time. Each copy counts from 1 upward at its
    own pace. There is \emph{one} file on disk --- so what are the \emph{two}
    things the operating system is running?
  \end{block}
  \begin{itemize}
    \item The file on disk does not change while it runs; it is just bytes.
    \item Each run has its own counter value, its own position in the code,
      its own stack.
    \item The OS must remember all of this for each run separately.
  \end{itemize}
  \vspace{0.1cm}
  \begin{alertblock}{Running thread for today}
    We follow \texttt{./counter} launched twice --- \texttt{P1} and
    \texttt{P2} --- and ask what the kernel must record about each.
  \end{alertblock}
\end{frame}

\begin{frame}{Roadmap: Four Questions About a Running Program}
  \begin{enumerate}
    \item \key{What is a process?} --- a program in execution, not the program
      file.
    \item \key{What does a process look like in memory?} --- text, data, heap,
      and stack.
    \item \key{What states can a process be in?} --- new, ready, running,
      waiting, terminated.
    \item \key{How does the kernel remember a process?} --- the process control
      block (PCB).
  \end{enumerate}
  \vspace{0.2cm}
  \muted{Last lecture organized the kernel itself. Today we name the unit of
  work that the kernel manages; next lecture we watch it move between queues.}
\end{frame}

\transitionslide{A program is a recipe written down. A process is somebody cooking it --- with the oven on, ingredients half-mixed, and a finger on line 7.}

% =============================================================================
% ACT 1 --- PROGRAM VS. PROCESS
% =============================================================================
\section{Program vs.\ Process}

\begin{frame}{Passive vs.\ Active}
  \centering
  \renewcommand{\arraystretch}{1.35}
  \begin{tabular}{p{0.42\textwidth}p{0.42\textwidth}}
    \textbf{Program} & \textbf{Process} \\
    \hline
    A file on disk (an executable) & A program in execution \\
    \emph{Passive}: a list of instructions & \emph{Active}: has a program
    counter and live resources \\
    One copy, however many runs & One per run, each independent \\
    Exists until deleted & Created, runs, and terminates \\
  \end{tabular}
  \vspace{0.35cm}
  \begin{itemize}
    \item A program becomes a process when its executable is \key{loaded into
      memory}.
    \item Two processes may run the same program; they are still two separate
      execution sequences.
  \end{itemize}
  \muted{Silberschatz et al., Ch.~3 \S3.1, pp.~101--103.}
\end{frame}

\begin{frame}{What Exactly Is a Process?}
  \begin{block}{Definition}
    A \key{process} is a program in execution: the unit of work in a modern
    time-sharing system. It consists of the program code together with its
    current activity --- the value of the program counter, the contents of the
    processor's registers, and its memory.
  \end{block}
  \begin{itemize}
    \item Batch systems called the same thing a \key{job}; the textbook uses
      \emph{job} and \emph{process} almost interchangeably.
    \item Stallings lists equivalent views: a program in execution, an
      instance of a program running on a computer, the entity that can be
      assigned to and executed on a processor.
  \end{itemize}
  \muted{Silberschatz et al., Ch.~3 \S3.1, p.~101; Stallings, Ch.~3 \S3.1,
  p.~108.}
\end{frame}

\begin{frame}{Worked Example: \texttt{./counter} Launched Twice}
  \begin{enumerate}
    \item Terminal A runs \texttt{./counter}. The OS creates \texttt{P1}:
      loads the code, allocates a stack, sets the program counter to the entry
      point.
    \item Terminal B runs \texttt{./counter}. The OS creates \texttt{P2}
      from the \emph{same} file, with a \emph{separate} stack and data.
    \item After a few seconds, \texttt{P1} has printed 812 and \texttt{P2}
      has printed 305.
    \item The variable \texttt{count} has two values at once --- one in each
      process's data.
    \item Both program counters may point into the same loop, at different
      instructions.
  \end{enumerate}
  \vspace{0.1cm}
  \begin{exampleblock}{Conclusion}
    Same text section, different activity. The OS tracks \texttt{P1} and
    \texttt{P2} as two processes.
  \end{exampleblock}
\end{frame}

\begin{frame}{Socratic Check 1}
  \begin{block}{Question}
    You open three browser windows of the same browser program. Is that one
    process or three? What would you need to know to decide?
  \end{block}
  \vspace{0.25cm}
  \begin{itemize}
    \item The number of \emph{windows} does not decide it; the number of
      \emph{execution sequences} the OS tracks does.
    \item One process can draw many windows; a program can also start several
      processes.
    \item The question to ask the OS is: how many PCBs exist for this program?
  \end{itemize}
\end{frame}

% =============================================================================
% ACT 2 --- PROCESS IN MEMORY
% =============================================================================
\section{A Process in Memory}

\transitionslide{Before the kernel can record a process, the process needs a place to live. Its address space has four regions, each with a job.}

% Coordinate map (P2): column x in [0,3.2]; regions bottom->top:
%   text  y[0,0.8]   data y[0.8,1.6]   heap y[1.6,2.4]
%   free  y[2.4,3.6] (dashed, arrows: heap grows up from 2.4, stack grows
%   down from 3.6)   stack y[3.6,4.4].
%   Address labels at x=-0.25 (anchor east): "max" at y=4.4, "0" at y=0.
%   Region notes at x=3.55 (anchor west), same y as each region's centre.
%   Arrows sit inside the free band at x=0.8 (up) and x=2.4 (down);
%   no label is placed on an arrow.
\begin{frame}{The Four Regions of an Address Space}
  \begin{columns}[c]
    \begin{column}{0.54\textwidth}
      \centering
      \begin{tikzpicture}[every node/.style={font=\small}]
        \draw[draw=primary-blue, fill=light-bg] (0,3.6) rectangle (3.2,4.4);
        \node at (1.6,4.0) {stack};
        \draw[draw=neutral, dashed] (0,2.4) rectangle (3.2,3.6);
        \draw[draw=primary-blue, fill=light-bg] (0,1.6) rectangle (3.2,2.4);
        \node at (1.6,2.0) {heap};
        \draw[draw=primary-blue, fill=light-bg] (0,0.8) rectangle (3.2,1.6);
        \node at (1.6,1.2) {data};
        \draw[draw=primary-blue, fill=light-bg] (0,0) rectangle (3.2,0.8);
        \node at (1.6,0.4) {text};
        \draw[observed-edge] (0.8,2.5) -- (0.8,3.0);
        \draw[observed-edge] (2.4,3.5) -- (2.4,3.0);
        \node[anchor=east, text=neutral] at (-0.25,4.4) {max};
        \node[anchor=east, text=neutral] at (-0.25,0) {0};
        \node[anchor=west, text=primary-gold, font=\footnotesize] at (3.55,4.0) {locals, returns};
        \node[anchor=west, text=primary-gold, font=\footnotesize] at (3.55,3.0) {free};
        \node[anchor=west, text=primary-gold, font=\footnotesize] at (3.55,2.0) {\texttt{malloc}};
        \node[anchor=west, text=primary-gold, font=\footnotesize] at (3.55,1.2) {globals};
        \node[anchor=west, text=primary-gold, font=\footnotesize] at (3.55,0.4) {code};
      \end{tikzpicture}
    \end{column}
    \begin{column}{0.42\textwidth}
      \begin{itemize}
        \item \key{Text}: the program code.
        \item \key{Data}: global variables.
        \item \key{Heap}: memory allocated at run time; grows upward.
        \item \key{Stack}: function parameters, return addresses, local
          variables; grows downward.
      \end{itemize}
      \vspace{0.1cm}
      \muted{Silberschatz et al., Ch.~3 \S3.1, pp.~101--102.}
    \end{column}
  \end{columns}
\end{frame}

\begin{frame}[fragile]{Worked Example: Where Does Each Variable Live?}
  \small
  \begin{columns}[T]
    \begin{column}{0.46\textwidth}
\begin{verbatim}
int count = 0;  /* (a) */
int main(void) {
  int step = 1;  /* (b) */
  int *log =
    malloc(64);  /* (c) */
  while (1) {
    count += step; /* (d) */
    print(count);
  }
}
\end{verbatim}
    \end{column}
    \begin{column}{0.50\textwidth}
      \begin{enumerate}
        \item[(a)] \texttt{count} is global --- \key{data}.
        \item[(b)] \texttt{step} is a local of \texttt{main} --- \key{stack}.
        \item[(c)] the 64 bytes \texttt{log} points to --- \key{heap};
          the pointer \texttt{log} itself --- \key{stack}.
        \item[(d)] the instructions of the loop --- \key{text}.
      \end{enumerate}
      \vspace{0.1cm}
      \begin{exampleblock}{For \texttt{P1} and \texttt{P2}}
        Same text; separate data, heap, and stack. \texttt{count} is 812 in
        one and 305 in the other.
      \end{exampleblock}
    \end{column}
  \end{columns}
\end{frame}

\begin{frame}{Process Image (Stallings' View)}
  \begin{itemize}
    \item Stallings calls the whole collection a \key{process image}:
    \begin{itemize}
      \item user program (text),
      \item user data (data and heap),
      \item stack,
      \item and the process's \key{control block}, the kernel's record of it.
    \end{itemize}
    \item The first three are what the process uses. The last is what the
      \emph{kernel} uses to manage it.
  \end{itemize}
  \vspace{0.2cm}
  \begin{block}{Key point}
    The memory regions describe what the process \emph{is doing}; the control
    block describes what the kernel \emph{knows about it}.
  \end{block}
  \muted{Stallings, Ch.~3 \S3.3, p.~126.}
\end{frame}

% =============================================================================
% ACT 3 --- PROCESS STATES
% =============================================================================
\section{Process States}

\transitionslide{A process is not always running. Most of the time it is waiting --- for the CPU or for the outside world --- and the kernel must know which.}

\begin{frame}{The Five States}
  \centering
  \footnotesize
  \renewcommand{\arraystretch}{1.3}
  \begin{tabular}{p{0.18\textwidth}p{0.66\textwidth}}
    \textbf{State} & \textbf{Meaning} \\
    \hline
    \key{new} & The process is being created. \\
    \key{ready} & The process is waiting to be assigned to a processor. \\
    \key{running} & Instructions are being executed. \\
    \key{waiting} & The process is waiting for some event (I/O completion,
    a signal). Stallings: \emph{blocked}. \\
    \key{terminated} & The process has finished execution. Stallings:
    \emph{exit}. \\
  \end{tabular}
  \vspace{0.3cm}
  \begin{alertblock}{The counting rule}
    On one processor, \textbf{at most one} process is running at any instant;
    many may be ready or waiting.
  \end{alertblock}
  \muted{Silberschatz et al., Ch.~3 \S3.1, p.~103; Stallings, Ch.~3 \S3.2,
  p.~111.}
\end{frame}

% Coordinate map (P2): nodes 2.3cm x 0.8cm.
%   new (-5.2,1.6)  ready (-2.0,0)  running (2.0,0)  terminated (5.2,1.6)
%   waiting (0,-2.2)
% Edges: new->ready (label above-left of midpoint, sloped);
%   ready->running bend left (arc above, label above arc);
%   running->ready bend left (arc below, label below arc);
%   running->waiting from running.south to waiting.east (label right);
%   waiting->ready from waiting.west to ready.south (label left);
%   running->terminated (label above, sloped).
% The two ready/running arcs leave the boxes at their east/west faces and stay
% between x=-0.85 and 0.85; no edge crosses a box.
\begin{frame}{The Process-State Diagram}
  \centering
  \begin{tikzpicture}[every node/.style={font=\small},
      st/.style={dag-node, minimum width=2.3cm, minimum height=0.8cm}]
    \node[st] (new) at (-5.2,1.6) {new};
    \node[st] (ready) at (-2.0,0) {ready};
    \node[st] (run) at (2.0,0) {running};
    \node[st] (term) at (5.2,1.6) {terminated};
    \node[st] (wait) at (0,-2.2) {waiting};
    \draw[observed-edge] (new) -- node[above, sloped, font=\footnotesize] {admitted} (ready);
    \draw[observed-edge] (ready.north east) to[bend left=30]
      node[above, font=\footnotesize] {scheduler dispatch} (run.north west);
    \draw[observed-edge] (run.south west) to[bend left=30]
      node[below, font=\footnotesize] {interrupt} (ready.south east);
    \draw[observed-edge] (run.south) |- node[pos=0.25, right, font=\footnotesize] {I/O or event wait} (wait.east);
    \draw[observed-edge] (wait.west) -| node[pos=0.75, left, font=\footnotesize] {I/O or event completion} (ready.south);
    \draw[observed-edge] (run) -- node[above, sloped, font=\footnotesize] {exit} (term);
  \end{tikzpicture}
  \vspace{0.15cm}
  \muted{After Silberschatz et al., Ch.~3 \S3.1, p.~103. Six transitions; each
  is caused by a specific event.}
\end{frame}

\begin{frame}{Reading the Diagram: What Causes Each Arrow?}
  \footnotesize
  \renewcommand{\arraystretch}{1.25}
  \begin{tabular}{p{0.27\textwidth}p{0.62\textwidth}}
    \textbf{Transition} & \textbf{Caused by} \\
    \hline
    new $\to$ ready & the OS admits the process (resources allocated) \\
    ready $\to$ running & the scheduler dispatches it to the CPU \\
    running $\to$ ready & an interrupt, e.g.\ the timer: its time is up \\
    running $\to$ waiting & it requests I/O or waits for an event (a blocking
    system call such as \texttt{read}) \\
    waiting $\to$ ready & the I/O completes or the event occurs \\
    running $\to$ terminated & it calls \texttt{exit} (or is killed) \\
  \end{tabular}
  \vspace{0.3cm}
  \begin{block}{Arrows that do not exist}
    waiting $\to$ running and ready $\to$ waiting. A waiting process must
    first become ready; a ready process is not executing, so it cannot request
    I/O.
  \end{block}
\end{frame}

\begin{frame}{Worked Trace: \texttt{P1} Reads a Key}
  \begin{enumerate}
    \item \texttt{P1} is created: \key{new}.
    \item The OS admits it: \key{ready}.
    \item The scheduler dispatches it: \key{running}; it increments
      \texttt{count}.
    \item The timer interrupt fires: back to \key{ready}; \texttt{P2} runs.
    \item \texttt{P1} is dispatched again: \key{running}. It calls
      \texttt{read} on the keyboard: \key{waiting}.
    \item The user presses a key; the interrupt completes the I/O:
      \key{ready}.
    \item \texttt{P1} is dispatched: \key{running}. It calls \texttt{exit}:
      \key{terminated}.
  \end{enumerate}
  \vspace{0.1cm}
  \muted{Seven states visited, six transitions taken; the blocking
  \texttt{read} is the system call from Lecture~3 seen from the process's side.}
\end{frame}

\begin{frame}{Worked Example: Counting States on One CPU}
  \begin{itemize}
    \item Five processes exist. \texttt{P1} is running; \texttt{P2} and
      \texttt{P3} wait for the disk; \texttt{P4} and \texttt{P5} are ready.
    \item \texttt{P1}'s time slice expires.
  \end{itemize}
  \vspace{0.15cm}
  \begin{exampleblock}{After the interrupt and one dispatch}
    \texttt{P1} $\to$ ready. The scheduler picks \texttt{P4} $\to$ running.\\
    Running: 1 (\texttt{P4}) \quad Ready: 2 (\texttt{P5}, \texttt{P1}) \quad
    Waiting: 2 (\texttt{P2}, \texttt{P3}).
  \end{exampleblock}
  \vspace{0.1cm}
  \begin{itemize}
    \item The total is still five: transitions move processes, they do not
      create them.
    \item With $k$ CPUs, at most $k$ processes run; the rest are ready or
      waiting.
  \end{itemize}
\end{frame}

\begin{frame}{Two-State and Five-State Models}
  \begin{itemize}
    \item Stallings starts from a \key{two-state} model: \emph{running} and
      \emph{not running}.
    \item It is too coarse: a not-running process might be ready to go, or
      blocked on I/O --- dispatching a blocked process wastes the CPU.
    \item Splitting \emph{not running} into \emph{ready} and \emph{blocked},
      and adding \emph{new} and \emph{exit}, gives the \key{five-state} model.
    \item Stallings later adds \emph{suspended} states for swapped-out
      processes (a seven-state model); we return to these with the
      medium-term scheduler.
  \end{itemize}
  \vspace{0.2cm}
  \muted{Stallings, Ch.~3 \S3.2, p.~111.}
\end{frame}

\begin{frame}{Socratic Check 2}
  \begin{block}{Question}
    A process in the waiting state receives the I/O it asked for. Why does it
    go to \emph{ready} instead of straight to \emph{running}?
  \end{block}
  \vspace{0.25cm}
  \begin{itemize}
    \item Another process is probably running on the CPU right now.
    \item Only the scheduler decides who runs next; the I/O completion only
      makes the process \emph{eligible}.
    \item Keeping ``eligible'' and ``chosen'' separate is what lets the OS
      apply a scheduling policy at all.
  \end{itemize}
\end{frame}

% =============================================================================
% ACT 4 --- THE PCB
% =============================================================================
\section{The Process Control Block}

\transitionslide{When P1 stops running, everything it was doing must be written down somewhere --- or it can never be resumed. That somewhere is the PCB.}

\begin{frame}{What Exactly Is a PCB?}
  \begin{block}{Definition}
    The \key{process control block} (PCB), also called a \key{task control
    block}, is the kernel data structure that represents one process. It is
    the repository for any information that may vary from process to process.
  \end{block}
  \begin{itemize}
    \item One PCB per process: \texttt{P1} and \texttt{P2} have two PCBs,
      even though they share one program file.
    \item The PCB lives in \emph{kernel} memory; the process cannot read or
      write it directly.
    \item Creating a process means creating a PCB; terminating it means
      eventually freeing the PCB.
  \end{itemize}
  \muted{Silberschatz et al., Ch.~3 \S3.1, p.~103.}
\end{frame}

% Coordinate map (P2): seven field boxes, each 4.6cm x 0.6cm, centred at x=0,
%   stacked downward from y=0 in steps of 0.6 (y = 0, -0.6, ..., -3.6), so
%   adjacent boxes share an edge and nothing overlaps. Title "PCB of P1"
%   at (0,0.6). No arrows; no labels outside the boxes.
\begin{frame}{Inside the PCB}
  \begin{columns}[c]
    \begin{column}{0.42\textwidth}
      \centering
      \begin{tikzpicture}[every node/.style={font=\small},
          fld/.style={draw=primary-blue, fill=light-bg, minimum width=4.6cm,
            minimum height=0.6cm, inner sep=1pt}]
        \node[text=primary-gold, font=\small\bfseries] at (0,0.6) {PCB of \texttt{P1}};
        \node[fld] at (0,0) {process state};
        \node[fld] at (0,-0.6) {process number (PID)};
        \node[fld] at (0,-1.2) {program counter};
        \node[fld] at (0,-1.8) {registers};
        \node[fld] at (0,-2.4) {memory limits};
        \node[fld] at (0,-3.0) {list of open files};
        \node[fld] at (0,-3.6) {$\cdots$};
      \end{tikzpicture}
    \end{column}
    \begin{column}{0.54\textwidth}
      \begin{itemize}
        \item \key{State}: new, ready, running, waiting, terminated.
        \item \key{PID}: the unique process number.
        \item \key{Program counter}: the next instruction to execute.
        \item \key{Registers}: accumulators, index registers, stack pointer,
          general-purpose registers, condition codes.
        \item \key{Memory limits}: where its regions are.
        \item \key{Open files}: what it is using.
      \end{itemize}
      \muted{After Silberschatz et al., Ch.~3 \S3.1, p.~103.}
    \end{column}
  \end{columns}
\end{frame}

\begin{frame}{The Seven Categories of PCB Information}
  \footnotesize
  \renewcommand{\arraystretch}{1.25}
  \begin{tabular}{p{0.27\textwidth}p{0.62\textwidth}}
    \textbf{Category} & \textbf{Examples} \\
    \hline
    Process state & new, ready, running, waiting, halted \\
    Program counter & address of the next instruction \\
    CPU registers & accumulators, index registers, stack pointers,
    general-purpose registers, condition codes \\
    CPU-scheduling information & priority, pointers to scheduling queues \\
    Memory-management information & base and limit registers, page tables or
    segment tables \\
    Accounting information & CPU time used, time limits, job or process
    numbers \\
    I/O status information & I/O devices allocated, list of open files \\
  \end{tabular}
  \vspace{0.25cm}
  \muted{Silberschatz et al., Ch.~3 \S3.1, pp.~103--104.}
\end{frame}

\begin{frame}{Why Save the PC and Registers?}
  \begin{itemize}
    \item When an interrupt takes the CPU away from \texttt{P1}, the CPU's
      registers still hold \texttt{P1}'s values.
    \item \texttt{P2} is about to overwrite them.
    \item So the kernel copies \texttt{P1}'s PC and registers \emph{into
      \texttt{P1}'s PCB} first.
    \item Later, it copies them back, and \texttt{P1} continues as if nothing
      happened.
  \end{itemize}
  \vspace{0.2cm}
  \begin{alertblock}{The PCB is what makes interruption safe}
    State information must be saved when an interrupt occurs, to allow the
    process to be continued correctly afterward. Saving one and loading another
    is the \key{context switch} --- the topic of the next lecture.
  \end{alertblock}
\end{frame}

\begin{frame}{Worked Example: Two PCBs, One Program}
  \centering
  \footnotesize
  \renewcommand{\arraystretch}{1.25}
  \begin{tabular}{p{0.28\textwidth}p{0.28\textwidth}p{0.28\textwidth}}
    \textbf{Field} & \textbf{PCB of \texttt{P1}} & \textbf{PCB of \texttt{P2}} \\
    \hline
    state & waiting (keyboard) & running \\
    PID & 4120 & 4127 \\
    program counter & inside \texttt{read} wrapper & inside the loop \\
    \texttt{count} (in data) & 812 & 305 \\
    open files & terminal A & terminal B \\
    text section & \multicolumn{2}{c}{shared: same \texttt{./counter} code} \\
  \end{tabular}
  \vspace{0.3cm}
  \begin{exampleblock}{Reading the table}
    Every row that differs is per-process information --- exactly what the PCB
    exists to hold. The PIDs are illustrative; the OS picks them.
  \end{exampleblock}
\end{frame}

\begin{frame}{Stallings' Three-Part View of the PCB}
  \begin{columns}[T]
    \begin{column}{0.32\textwidth}
      \begin{block}{Identification}
        PID, parent's PID, user ID.
      \end{block}
    \end{column}
    \begin{column}{0.32\textwidth}
      \begin{block}{Processor state}
        User-visible registers, control and status registers (PC, flags),
        stack pointers.
      \end{block}
    \end{column}
    \begin{column}{0.32\textwidth}
      \begin{block}{Process control}
        Scheduling and state, links to other processes, privileges, memory and
        resource ownership.
      \end{block}
    \end{column}
  \end{columns}
  \vspace{0.25cm}
  \begin{itemize}
    \item Same content as Silberschatz's seven categories, grouped by
      \emph{purpose}.
    \item Processor-state information is what a context switch saves and
      restores.
  \end{itemize}
  \muted{Stallings, Ch.~3 \S3.3, p.~126.}
\end{frame}

\begin{frame}{A Real PCB: Linux \texttt{task\_struct}}
  \begin{itemize}
    \item In Linux the PCB is the \key{process descriptor}, a C structure
      of type \texttt{task\_struct}.
    \item It holds the state, the PID, scheduling information, a pointer to
      the memory descriptor, the open-file table, parent and child links,
      and much more.
    \item Its \texttt{state} field records values such as
      \texttt{TASK\_RUNNING}, \texttt{TASK\_INTERRUPTIBLE}, and
      \texttt{TASK\_UNINTERRUPTIBLE}.
    \item The kernel links all descriptors into a list so it can walk every
      process.
  \end{itemize}
  \vspace{0.2cm}
  \begin{block}{Note the naming}
    Linux uses \texttt{TASK\_RUNNING} for both ready and running: the
    textbook's two states map to one flag plus ``is it on a CPU right now?''
  \end{block}
  \muted{Bovet \& Cesati, Ch.~3, pp.~79--130.}
\end{frame}

\begin{frame}{Socratic Check 3}
  \begin{block}{Question}
    Why can the PCB not be stored in the process's own stack or heap, where it
    would be cheaper to reach?
  \end{block}
  \vspace{0.25cm}
  \begin{itemize}
    \item The process could then overwrite its own state, PID, or memory
      limits --- and break protection.
    \item The kernel must read the PCB even when the process is not running
      and its memory may be swapped out.
    \item So PCBs live in kernel memory, reachable only in kernel mode.
  \end{itemize}
\end{frame}

% =============================================================================
% ACT 5 --- PUTTING IT TOGETHER
% =============================================================================
\section{Putting It Together}

\transitionslide{A process is its address space plus its PCB. One more refinement --- threads --- and then the misconceptions.}

\begin{frame}{Threads: A Preview}
  \begin{itemize}
    \item The model so far: a process performs a \emph{single thread of
      execution} --- one program counter.
    \item Modern systems let a process have \key{multiple threads}, each with
      its own PC and registers, sharing the same address space.
    \item Then the PCB is extended with per-thread information.
  \end{itemize}
  \vspace{0.2cm}
  \begin{exampleblock}{Example}
    A word processor with one thread cannot type characters and run the spell
    checker at the same time; with two threads it can.
  \end{exampleblock}
  \muted{Silberschatz et al., Ch.~3 \S3.1, pp.~104--105; threads are treated
  fully in Ch.~4.}
\end{frame}

\begin{frame}{Common Misconceptions}
  \begin{columns}[T]
    \begin{column}{0.48\textwidth}
      \begin{block}{Wrong}
        \begin{itemize}
          \item ``A process is the program file.''
          \item ``Two runs of one program share one PCB.''
          \item ``I/O completion puts a process back on the CPU.''
          \item ``Many processes can be running on one CPU.''
        \end{itemize}
      \end{block}
    \end{column}
    \begin{column}{0.48\textwidth}
      \begin{block}{Right}
        \begin{itemize}
          \item A process is a program \emph{in execution}.
          \item One PCB per process, even for the same program.
          \item It goes to \emph{ready}; the scheduler decides.
          \item At most one at a time per CPU; others are ready or waiting.
        \end{itemize}
      \end{block}
    \end{column}
  \end{columns}
\end{frame}

\begin{frame}{Takeaways}
  \begin{itemize}
    \item A \key{process} is a program in execution; the program is passive,
      the process is active.
    \item Its address space has \key{text}, \key{data}, \key{heap}, and
      \key{stack}; two processes running one program share text only.
    \item It moves among five \key{states}: new, ready, running, waiting,
      terminated --- six legal transitions, each caused by an event.
    \item The \key{PCB} records everything per-process: state, PC,
      registers, scheduling, memory, accounting, I/O status.
    \item Saving the PC and registers into the PCB is what makes interruption
      and resumption possible.
  \end{itemize}
\end{frame}

\begin{frame}{Next: Process States, Schedulers, and Context Switch}
  \begin{itemize}
    \item Today: what a process is, how it is laid out, which states it moves
      through, and how the PCB records it.
    \item Next lecture: the queues that hold PCBs, the long-, short-, and
      medium-term schedulers that move them, and the context switch that saves
      one PCB and loads another.
  \end{itemize}
  \vspace{0.3cm}
  \begin{block}{Reading}
    Silberschatz et al., \emph{Operating System Concepts}, Ch.~3 \S3.2
    \cite{SilberschatzOSC_operating_system_concepts}; Stallings,
    \emph{Operating Systems: Internals and Design Principles}, Ch.~3 \S3.2 and
    \S3.4 \cite{StallingsOS_internals_design_principles}; Bovet \& Cesati,
    Ch.~3 \cite{BovetCesati_understanding_linux_kernel}.
  \end{block}
\end{frame}

\begin{frame}{Further Reading (Page-Cited)}
  \footnotesize
  \begin{itemize}
    \item A.~Silberschatz, P.~B.~Galvin, G.~Gagne, \emph{Operating System
      Concepts}, 8e --- Ch.~3 \S3.1 (the process, p.~101; process state,
      p.~103; process control block, p.~103; threads, pp.~104--105).
    \item W.~Stallings, \emph{Operating Systems: Internals and Design
      Principles}, 6e --- Ch.~3 \S3.1 (what is a process, p.~108), \S3.2
      (process states, p.~111), \S3.3 (process description, p.~126).
    \item D.~P.~Bovet \& M.~Cesati, \emph{Understanding the Linux Kernel}, 3e
      --- Ch.~3 (processes and the process descriptor, pp.~79--130), for the
      concrete Linux \texttt{task\_struct}.
  \end{itemize}
  \vspace{0.2cm}
  \muted{Page-verified via each book's \texttt{index.md} under
  \texttt{master\_supporting\_docs/CS502/supporting\_books/}.}
\end{frame}

\begin{frame}[allowframebreaks]{References}
  \footnotesize
  \bibliographystyle{plain}
  \bibliography{../../Bibliography_base}
\end{frame}

\end{document}
